\begin{abstract}
\textbf{[Placeholder abstract --- to be written once the report is finalised.]} This project investigates whether decomposing a long-horizon robotic manipulation task into language-conditioned subtasks, sequenced by an orchestrator that monitors outcomes and recovers from failures at two levels --- retrying the same sub-goal and redirecting to an alternative orange to preserve task progress --- improves robustness over running a single policy end-to-end. The study is carried out on a simulated pick-and-place task and compares teleoperated and autonomously augmented training data across monotask and orchestrated subtask formulations. A full summary of the method and findings will be added here.

\end{abstract}

\begin{figure}[H]
    \centering
    \resizebox{\textwidth}{!}{%
    \begin{tikzpicture}[
        font=\small, >=Latex, node distance=10mm,
        stage/.style={draw, thick, rounded corners=3pt, align=center,
                      minimum height=12mm, inner sep=4pt},
        data/.style={stage, draw=teal!60!black, fill=teal!10},
        model/.style={stage, draw=orange!85!black, fill=orange!16},
        orchs/.style={stage, draw=blue!60!black, fill=blue!10},
        env/.style={stage, draw=black!55, fill=black!7},
        srcbox/.style={draw, rounded corners=2pt, draw=teal!55!black, fill=teal!5,
                       align=center, font=\footnotesize, minimum height=8mm, text width=20mm, inner sep=2pt},
        group/.style={draw, dashed, rounded corners=4pt, inner sep=3.4mm},
        arr/.style={->, thick, draw=black!70},
        loop/.style={->, thick, draw=teal!55!black, dashed},
        gamepad/.pic={
            \draw[fill=teal!22, draw=teal!65!black, line width=0.5pt, line join=round]
                (-0.5,0.08)
                .. controls (-0.5,0.24) and (-0.34,0.26) .. (-0.18,0.26)
                -- (0.18,0.26)
                .. controls (0.34,0.26) and (0.5,0.24) .. (0.5,0.08)
                -- (0.5,-0.10)
                .. controls (0.5,-0.34) and (0.30,-0.36) .. (0.18,-0.22)
                .. controls (0.07,-0.10) and (-0.07,-0.10) .. (-0.18,-0.22)
                .. controls (-0.30,-0.36) and (-0.5,-0.34) .. (-0.5,-0.10)
                -- cycle;
            \fill[white] (-0.34,-0.03) rectangle (-0.18,0.03);
            \fill[white] (-0.29,-0.08) rectangle (-0.23,0.08);
            \fill[white] (0.26,0.07) circle (0.025);
            \fill[white] (0.26,-0.07) circle (0.025);
            \fill[white] (0.19,0.0) circle (0.025);
            \fill[white] (0.33,0.0) circle (0.025);
        },
    ]
        % --- LeRobot group: datasets feed fine-tuning ---
        \node[data] (data) {\textbf{Datasets}\\\footnotesize Teleop\,$+$\,Auto};
        \node[model, right=of data] (train) {\textbf{SmolVLA}\\\footnotesize fine-tuning};

        % --- Isaac Sim group: orchestrator prompts the policy; policy drives the robot ---
        \node[model, right=26mm of train] (vla) {\textbf{SmolVLA}\\\footnotesize policy};
        \node[env, right=16mm of vla] (robot) {\textbf{Simulated}\\\textbf{robot}};
        \coordinate (vrmid) at ($(vla)!0.5!(robot)$);
        \node[orchs, above=9mm of vrmid] (orch) {\textbf{Orchestrator}\\\footnotesize GRASP / LIFT / PLACE};

        % --- data sources (left of the dataset stage) ---
        \node[srcbox] (base) [left=13mm of data, yshift=-5mm] {LightwheelAI baseline};
        \node[srcbox] (tele) [above=4mm of base] {Teleoperation};
        \pic[scale=0.72] at ([yshift=6mm]tele.north) {gamepad};

        % --- framework / environment grouping ---
        \begin{scope}[on background layer]
            \node[group, draw=black!45, fill=black!3, fit=(data)(train),
                  label={[font=\footnotesize\bfseries, black!55]above:LeRobot}] {};
            \node[group, draw=green!45!black, fill=green!5, fit=(vla)(robot)(orch),
                  label={[font=\footnotesize\bfseries, green!45!black]above:Isaac Sim --- evaluation / inference}] (isim) {};
        \end{scope}

        % --- arrows ---
        \draw[arr] (tele.east) -- (data.west);
        \draw[arr] (base.east) -- (data.west);
        \draw[arr] (data) -- (train);
        \draw[arr] (train.east) -- node[above, font=\scriptsize, align=center]{fine-tuned\\policy} (vla.west);

        % orchestrator -> policy -> robot -> orchestrator (inside Isaac Sim)
        \draw[arr] (orch) -- node[midway, fill=white, inner sep=0.6pt, font=\scriptsize]{prompt} (vla);
        \draw[arr] (vla) -- node[midway, below, font=\scriptsize]{actions} (robot);
        \draw[arr] (robot) -- node[midway, fill=white, inner sep=0.6pt, font=\scriptsize]{state} (orch);

        % --- autonomous-data feedback loop ---
        \draw[loop] (isim.south) -- ++(0,-0.7) -| (data.south);
        \node[font=\footnotesize, text=teal!45!black, align=center]
            at ($(data.south)!0.5!(isim.south)+(0,-0.35)$)
            {successful subtask rollouts \;(autonomous \emph{Auto} data)};
    \end{tikzpicture}%
    }
    \caption{Overview of the project pipeline.}
    \label{fig:pipeline}
\end{figure}
